\documentclass[11pt,a4paper]{article}

\usepackage[T1]{fontenc}
\usepackage[utf8]{inputenc}
\usepackage{lmodern}
\usepackage{microtype}
\usepackage[margin=2.3cm,top=2.6cm,bottom=2.4cm]{geometry}
\usepackage{booktabs}
\usepackage{tabularx}
\usepackage{array}
\usepackage{longtable}
\usepackage{pdflscape}
\usepackage[table,dvipsnames]{xcolor}
\usepackage{titlesec}
\usepackage{enumitem}
\usepackage{fancyhdr}
\usepackage{caption}
\usepackage{amssymb}
\usepackage{graphicx}
\usepackage[colorlinks=true,linkcolor=black,urlcolor=black,citecolor=black]{hyperref}

% ---- palette -----------------------------------------------------------
% Monochrome: hierarchy comes from weight, size and rules -- not colour.
\definecolor{accent}{HTML}{000000}
\definecolor{accentlt}{HTML}{FFFFFF}
\definecolor{rule}{HTML}{000000}
\definecolor{warn}{HTML}{000000}
\definecolor{warnlt}{HTML}{FFFFFF}
\definecolor{ink}{HTML}{000000}
\definecolor{muted}{HTML}{000000}

\color{ink}

% ---- headings ----------------------------------------------------------
\titleformat{\section}
  {\normalfont\Large\bfseries}{\thesection}{0.7em}{}[\vspace{-0.55em}{\color{black}\rule{\linewidth}{0.7pt}}]
\titlespacing*{\section}{0pt}{1.7em}{0.7em}

\titleformat{\subsection}
  {\normalfont\large\bfseries\color{ink}}{\thesubsection}{0.6em}{}
\titlespacing*{\subsection}{0pt}{1.3em}{0.5em}

\titleformat{\subsubsection}
  {\normalfont\normalsize\bfseries\color{muted}}{\thesubsubsection}{0.6em}{}

% ---- header / footer ---------------------------------------------------
\pagestyle{fancy}
\setlength{\headheight}{22pt}
\addtolength{\topmargin}{-10pt}
\fancyhf{}
\renewcommand{\headrulewidth}{0.3pt}
\renewcommand{\headrule}{\hbox to\headwidth{\color{rule}\leaders\hrule height \headrulewidth\hfill}}
\fancyhead[L]{\small\color{muted}Submission 1 — Dataset, EDA and Related Work}
\fancyhead[R]{\small\color{muted}Cricket Shot Attention CNN}
\fancyfoot[C]{\small\color{muted}\thepage}

% ---- helpers -----------------------------------------------------------
\newcommand{\blank}{\rule[-0.2ex]{1.8em}{1.1pt}}

\newenvironment{keybox}
  {\par\medskip\noindent\begin{tabular}{@{}p{0.3em}@{\hspace{0.8em}}p{\dimexpr\linewidth-2.2em}@{}}
   {\color{accent}\vrule width 2.5pt} &}
  {\end{tabular}\par\medskip}

\newenvironment{warnbox}
  {\par\medskip\noindent\begin{tabular}{@{}p{0.3em}@{\hspace{0.8em}}p{\dimexpr\linewidth-2.2em}@{}}
   {\color{warn}\vrule width 2.5pt} &}
  {\end{tabular}\par\medskip}

\setlist[itemize]{leftmargin=1.3em,itemsep=0.25em,topsep=0.4em}
\setlist[enumerate]{leftmargin=1.5em,itemsep=0.3em,topsep=0.4em}

\renewcommand{\arraystretch}{1.25}
\captionsetup{font=small,labelfont={bf,color=accent},skip=6pt}

\setlength{\parskip}{0.55em}
\setlength{\parindent}{0pt}

% ========================================================================
\begin{document}

% ---- title block -------------------------------------------------------
\begin{center}
{\color{muted}\small CNN PROJECT \textbullet\ SUBMISSION 1 OF 4}\\[0.9em]
{\LARGE\bfseries\color{accent} Attention-Guided Fine-Grained Classification\\[0.25em]
of Cricket Batting Shots from Single Frames}\\[1.1em]
{\large Dataset, Exploratory Analysis and Related Work}\\[1.4em]
{\color{muted}\small
Repository: \href{https://github.com/Rayhan-who/cricket-shot-attention-cnn}{github.com/Rayhan-who/cricket-shot-attention-cnn}}
\end{center}

\vspace{0.6em}
{\hrule height 0.8pt}
\vspace{1.2em}

% ========================================================================
\section*{Summary}
\addcontentsline{toc}{section}{Summary}

This submission establishes the dataset and evaluation protocol for a four-part CNN
project on cricket batting shot recognition. The task is \textbf{multi-class,
single-label image classification} over four shot types, performed on
\textbf{individual still frames} rather than video sequences.

Three outcomes are delivered here. First, a complete inventory and exploratory analysis
of the dataset. Second, a \textbf{duplicate audit} establishing that the published
dataset contains offline-augmented copies of the same source photographs, and a
correspondingly leakage-safe train/validation/test split that all later submissions
reuse unchanged. Third, a survey of prior work establishing that the field is almost
entirely video-based, that published accuracies in it are known to be difficult to
reproduce, and that the contribution of attention has not previously been isolated.

\begin{keybox}
\textbf{Framing.} Almost all prior cricket shot classification work models temporal
sequences with LSTM or GRU layers over multiple frames. This project removes that signal
entirely and asks how much of the task is recoverable from a single frame --- and
whether spatial attention can close part of the resulting gap.
\end{keybox}

% ========================================================================
\section{Dataset information}

\subsection{Provenance}
\label{sec:prov}

\begin{keybox}
\textbf{All provenance fields verified against the dataset page} on 7 November 2026.
Kaggle reports the licence as \emph{Unknown}. The current release is version~2, last
updated on 20 September 2022 according to Kaggle's dataset metadata.
\end{keybox}

\begin{tabularx}{\linewidth}{@{}>{\bfseries}p{4.2cm}X@{}}
\toprule
Field & Value \\
\midrule
Name & Cricket Shots Dataset \\
Source & \href{https://www.kaggle.com/datasets/aneesh10/cricket-shot-dataset}{kaggle.com/datasets/aneesh10/cricket-shot-dataset} \\
Author & Aneesh Aparajit (Kaggle user \texttt{aneesh10}) \\
Licence & \textbf{Unknown} --- as stated on the Kaggle page; the author has specified no licence \\
Year published & 2022 --- version 2, last updated 20 September 2022 \\
Task type & Multi-class, single-label image classification \\
Update frequency & Never --- the dataset is static, so these measurements stay reproducible \\
Community signal & Usability 7.50, 56 upvotes \\
Number of classes & 4 \\
Class names & \texttt{drive}, \texttt{legglance-flick}, \texttt{pullshot}, \texttt{sweep} \\
Construction & Images collected from the internet, then expanded by offline augmentation --- stated by the author on the dataset page \\
\bottomrule
\end{tabularx}

\begin{warnbox}
\textbf{The licence is Unknown, and that has consequences.} Kaggle displays
\emph{Unknown} when an uploader has specified no licence at all, which means no
permission to redistribute has been granted. Use of the images for coursework analysis
is defensible as academic fair dealing, but the dataset itself must not be re-uploaded,
mirrored, or committed to a public repository. This project therefore references the
dataset by link and commits only derived artefacts --- the split manifest, measured
statistics, and a small number of samples reproduced for scholarly illustration.
\end{warnbox}

\subsection{Class definitions}

\begin{tabularx}{\linewidth}{@{}>{\ttfamily}p{3.4cm}X@{}}
\toprule
\textnormal{\bfseries Class} & \textbf{Shots included} \\
\midrule
drive & Cover drive, straight drive, off drive \\
legglance-flick & Leg glance and flick \\
pullshot & Pull \\
sweep & Sweep \\
\bottomrule
\end{tabularx}

\subsection{Measured properties}

Computed by the accompanying notebook (\texttt{01\_eda\_and\_related\_work.ipynb}).

\begin{tabularx}{\linewidth}{@{}>{\bfseries}p{6.2cm}X@{}}
\toprule
Property & Measured value \\
\midrule
Total images & 4{,}724 \\
Images per class & \texttt{pullshot} 1{,}260 \textbullet\ \texttt{drive} 1{,}224 \textbullet\ \texttt{legglance-flick} 1{,}120 \textbullet\ \texttt{sweep} 1{,}120 \\
Width (min / median / max) & 168 / 275 / 1280 px \\
Height (min / median / max) & 157 / 191 / 1280 px \\
Distinct size combinations & 84 \\
Colour modes & \textbf{Mixed} --- 4{,}586 RGBA (4 channels), 138 RGB (3 channels) \\
File formats & \textbf{Mixed} --- PNG 4{,}586, JPEG 137, WebP 1 \\
Total size on disk & 643.8 MB \\
Corrupt or unreadable files & 0 \\
\bottomrule
\end{tabularx}

\begin{warnbox}
\textbf{Two properties here are not cosmetic.} First, 97\% of images are \textbf{RGBA},
carrying a fourth alpha channel. ImageNet-pretrained backbones expect three channels, so
every image must be explicitly converted to RGB before training; feeding RGBA tensors
through a three-channel stem either raises a shape error or, worse, silently
misinterprets the data. Second, the file formats are mixed, so a minority of images
carry JPEG compression artefacts that the PNG majority does not. Both are handled in
Section~\ref{sec:prep}.
\end{warnbox}

\subsection{Why this dataset matters in its domain}

Cricket has an audience exceeding one billion people, and the broadcast and coaching
workflows built around it remain substantially manual. Automatic batting-shot
recognition supports three applications that recur throughout the literature:

\begin{itemize}
\item \textbf{Highlight generation.} Broadcasters currently tag shot types by hand
  across many hours of footage. Automatic labelling is a direct cost reduction.
\item \textbf{Coaching and technique analysis.} Shot-type identification is the
  precondition for any downstream biomechanical feedback --- a system cannot assess the
  execution of a pull shot until it knows a pull shot was attempted.
\item \textbf{Automated commentary and context-aware advertising.} Both require
  machine-readable understanding of what is happening on the field.
\end{itemize}

What makes this a genuinely difficult classification problem, rather than a trivial one,
is that all four classes share an identical scene grammar: one batsman, one pitch, the
same protective equipment, comparable stadium backdrops, and frequently a crowd and
sightscreen in the background. The class label is determined almost entirely by
\textbf{body posture and bat angle}, which occupy a small fraction of the total frame
area. Prior work makes precisely this observation --- the visual similarity between
batting shots is what renders hand-crafted feature extraction impractical and motivates
learned representations~\cite{sen2021}.

This property is also what makes the dataset suitable for the later submissions. A
coarse or visually separable dataset would saturate near ceiling accuracy under any
reasonable pretrained backbone, leaving no measurable headroom for the attention module
introduced in Submission 4 to demonstrate an effect. Fine-grained difficulty is a
requirement of the experimental design, not an inconvenience.

% ========================================================================
\section{Exploratory data analysis}

\subsection{What was measured}

\begin{tabularx}{\linewidth}{@{}p{5.0cm}X@{}}
\toprule
\textbf{Analysis} & \textbf{Purpose} \\
\midrule
Class distribution & Determine whether accuracy alone is a valid headline metric, and whether class weighting is required \\
Sample grid, 5 per class & Visual assessment of background clutter, subject scale, and inter-class similarity \\
Dimension and aspect histograms & Decide the resize and crop strategy \\
Perceptual duplicate audit & Quantify redundancy and prevent train/test leakage \\
Grouped stratified split & Produce a frozen, reusable partition for Submissions 2--4 \\
\bottomrule
\end{tabularx}

\subsection{Sample images}

Figure~\ref{fig:samples} shows two representative frames from each class. They were
chosen to span the range of quality present in the dataset rather than to show the
clearest examples, because the variation is itself a finding.

\begin{figure}[htbp]
\centering
\newcommand{\simg}[1]{\raisebox{-0.5\height}{\includegraphics[width=0.405\linewidth]{#1}}}
\setlength{\tabcolsep}{3pt}
\begin{tabular}{@{}>{\bfseries\footnotesize\ttfamily}p{2.4cm}cc@{}}
drive &
  \simg{figures/drive_a.jpg} & \simg{figures/drive_b.jpg} \\[10pt]
legglance-flick &
  \simg{figures/leg_a.jpg} & \simg{figures/leg_b.jpg} \\[10pt]
pullshot &
  \simg{figures/pull_a.jpg} & \simg{figures/pull_b.jpg} \\[10pt]
sweep &
  \simg{figures/sweep_a.jpg} & \simg{figures/sweep_b.jpg} \\[4pt]
\end{tabular}
\caption{Two samples per class. Note the shared scene grammar across all four rows:
one batsman, one pitch, comparable kit and comparable backdrops. The class label is
carried almost entirely by body posture and bat angle. Note also the range of capture
quality --- daylight and floodlit matches, net practice and match play, sharp frames and
degraded broadcast captures with motion blur and colour cast.}
\label{fig:samples}
\end{figure}

Four observations follow directly from these samples and inform the preprocessing plan
in Section~\ref{sec:prep}.

\textbf{The subject occupies a minority of the frame.} In most samples the batsman
covers well under half the image area. The remainder is turf, crowd, sightscreen and
advertising hoardings --- none of which carries class information, and all of which
competes for the network's representational capacity. This is the specific condition
that spatial attention is intended to address in Submission 4.

\textbf{Backgrounds are correlated with capture context, not with class.} Advertising
boards, crowd density and floodlighting vary widely but are unrelated to which shot is
being played. A network that latches onto them will appear to perform well on the
training distribution and fail on new footage.

\textbf{Capture quality is highly variable.} The dataset mixes sharp high-resolution
match photography with low-resolution broadcast captures showing motion blur,
compression artefacts and strong colour casts. This argues for brightness and contrast
jitter during augmentation, and against any preprocessing step that assumes consistent
colour balance.

\textbf{Shot geometry is side-dependent.} Comparing the \texttt{sweep} and
\texttt{legglance-flick} rows makes clear that these classes are distinguished partly by
which side of the body the ball is played to, and by the degree of knee bend. This is
the direct justification for excluding horizontal flipping from the augmentation
pipeline --- see the note in Section~\ref{sec:prep}.

\subsection{Class balance}

\begin{tabularx}{\linewidth}{@{}>{\bfseries}p{6.2cm}X@{}}
\toprule
Quantity & Value \\
\midrule
Largest class & \texttt{pullshot} --- 1{,}260 images (26.67\%) \\
Smallest class & \texttt{legglance-flick} --- 1{,}120 images (23.71\%) \\
Imbalance ratio & 1.12 : 1 \\
Verdict & \textbf{Balanced} \\
\bottomrule
\end{tabularx}

At 1.12\,:\,1 the dataset is close to uniform. \textbf{Consequence:} plain accuracy is a
valid headline metric and no class weighting or resampling is required. Macro-averaged
F1 is still reported from Submission 2 onward, since per-class behaviour remains
informative even when the prior is flat, but it is not needed to correct a skew.

This is a clean negative finding and it simplifies the remaining submissions: one fewer
confounding variable between the Submission 3 baseline and the Submission 4 attention
model.

\subsection{A resolution confound}

The per-class geometry is \emph{not} uniform, and this matters more than the class counts
do:

\begin{tabularx}{\linewidth}{@{}>{\ttfamily}p{3.6cm}rrX@{}}
\toprule
\textnormal{\bfseries Class} & \textbf{Median w} & \textbf{Median h} & \textbf{Mean file size} \\
\midrule
drive           & 262 & 192 & 87.6 kB \\
legglance-flick & 410 & 388 & \textbf{302.5 kB} \\
pullshot        & 275 & 183 & 90.3 kB \\
sweep           & 278 & 181 & 88.9 kB \\
\bottomrule
\end{tabularx}

\texttt{legglance-flick} images are markedly larger than the other three classes --- a
median area roughly three times greater and a mean file size over three times larger.
The other three classes are tightly clustered.

\begin{warnbox}
\textbf{This is a potential shortcut feature.} A systematic difference in native
resolution that correlates with the label is exactly the kind of signal a CNN will
exploit in preference to learning posture. It suggests the \texttt{legglance-flick}
images were sourced or processed differently from the rest. Resizing every image to a
common $224 \times 224$ removes the raw dimension cue, but any residual artefact of
upsampling the smaller classes may persist. If \texttt{legglance-flick} shows anomalously
high per-class recall in Submission 2, this is the first thing to investigate.
\end{warnbox}

\subsection{Duplicate audit and leakage control}
\label{sec:dupaudit}

\begin{warnbox}
\textbf{This is the most consequential finding of Submission 1.} The dataset was
assembled by collecting images and then augmenting them offline before publication.
Multiple files are therefore transformations of a single source photograph. A
conventional per-file random split distributes those transformations across both the
training and test partitions, so the model is evaluated on images it has effectively
already seen, and \emph{every} accuracy figure reported downstream is inflated.
\end{warnbox}

\subsubsection{Method}

A 64-bit perceptual \emph{difference hash} (dHash) is computed for each image. Each
image is converted to greyscale, resized to $9 \times 8$ pixels, and the sign of the
horizontal gradient between adjacent pixels is recorded as one bit per comparison. The
resulting fingerprint is insensitive to minor changes in scale, compression and
brightness, but responds to genuine differences in content. Images sharing a hash are
collected into a \textbf{duplicate group}.

The implementation depends on nothing beyond Pillow and NumPy, so it runs unchanged on
Kaggle, Colab or a local machine.

\subsubsection{Findings}

\begin{tabularx}{\linewidth}{@{}>{\bfseries}p{7.0cm}X@{}}
\toprule
Quantity & Value \\
\midrule
Unique perceptual hashes & 4{,}116 of 4{,}724 images \\
Duplicate groups & 345 \\
Images inside a duplicate group & 953 \\
Redundant copies & 608 \quad (12.9\% of the dataset) \\
Largest single group & 10 copies of one image \\
Groups spanning more than one class & \textbf{0} \\
\textbf{Effective dataset size} & \textbf{4{,}116 independent images} \\
\bottomrule
\end{tabularx}

Two things follow.

\textbf{The labelling is sound.} No duplicate group spans more than one class, so there
are no cases of perceptually identical images carrying contradictory ground truth. The
dataset is careless about redundancy but not about annotation.

\textbf{The redundancy is material but moderate.} Roughly 20\% of files sit in a
multi-member group. Under a naive per-file split, approximately 17\% of the test set
would have a near-twin in training. Treating those as effectively memorised gives an
estimated accuracy inflation on the order of \textbf{two to three percentage points}
--- enough to change a model ranking, not enough to invalidate the task.

\begin{warnbox}
\textbf{12.9\% is a lower bound, not a measurement of total redundancy.} Exact dHash
collision detects only near-identical images. If the published augmentation included
geometric transformations --- rotation, flipping, cropping --- those copies produce
entirely different hashes and are \emph{not} counted here. A Hamming-distance variant,
treating hashes within a small bit radius as equivalent, would establish the true figure.
That extension is noted as future work.
\end{warnbox}

\subsubsection{Resulting split policy}

The train/validation/test partition is constructed over \emph{duplicate groups},
stratified by class, in a 70/15/15 ratio with seed 42. An assertion verifies that no
group's members appear in more than one partition; it passes.

\begin{tabularx}{\linewidth}{@{}>{\ttfamily}p{3.8cm}rrrX@{}}
\toprule
\textnormal{\bfseries Class} & \textbf{Train} & \textbf{Val} & \textbf{Test} & \textbf{Total} \\
\midrule
drive           & 862 & 180 & 182 & 1{,}224 \\
legglance-flick & 781 & 160 & 179 & 1{,}120 \\
pullshot        & 888 & 186 & 186 & 1{,}260 \\
sweep           & 790 & 156 & 174 & 1{,}120 \\
\midrule
\textnormal{\bfseries Total} & \textbf{3{,}321} & \textbf{682} & \textbf{721} & \textbf{4{,}724} \\
\bottomrule
\end{tabularx}

The result is written to \texttt{results-data/split\_manifest.csv} and committed to the
repository.

\begin{keybox}
\textbf{Submissions 2, 3 and 4 load this manifest. They do not re-split.} If the
partition changed between Submission 3 and Submission 4, the attention comparison would
be two unrelated training runs rather than a controlled experiment, and no difference
between them could be attributed to the attention module.
\end{keybox}

\subsection{Written findings}

\textbf{Finding 1 --- Balance.} The dataset contains 4{,}724 images across four classes.
The largest class, \texttt{pullshot}, holds 1{,}260 images; the smallest,
\texttt{legglance-flick}, holds 1{,}120. The imbalance ratio is 1.12\,:\,1, which is
effectively balanced. \emph{Consequence:} accuracy alone is a valid headline metric and
no class weighting is applied. Macro-averaged F1 is still reported from Submission 2
onward for per-class visibility, but not as a correction for skew.

\textbf{Finding 2 --- Redundancy and leakage risk.} Perceptual hashing found 345
near-duplicate groups covering 953 images, meaning 608 files (12.9\%) are redundant
copies of another image; the largest single group contains ten copies of one photograph.
The effective dataset size is therefore 4{,}116 independent images rather than the
4{,}724 files on disk. This follows directly from the dataset having been augmented
offline before publication --- a construction the author states explicitly on the
dataset page, which describes images gathered from the internet and then augmented
to produce the released collection. The redundancy is therefore an expected
property of the data rather than an anomaly, and the fact that it is documented by
the uploader makes its omission in downstream work harder to excuse. \emph{Consequence:} all partitioning is performed on
duplicate-group identity rather than file identity. A naive per-file random split would
place augmented copies of the same source photograph in both train and test; the
estimated inflation is two to three percentage points. No duplicate group spans more
than one class, so the annotation itself is sound.

\textbf{Finding 3 --- Class similarity and a resolution confound.} All four classes
share an identical scene grammar: one batsman on a pitch, comparable kit, similar
stadium backdrops, frequently a crowd and sightscreen behind. The discriminative signal
is confined to body posture and bat angle within a minority of the frame. Inspection of the
sample grid identifies \texttt{drive} and \texttt{legglance-flick} as the closest pair.
Both are played from an upright, predominantly side-on stance with the bat travelling in
a near-vertical arc; an on-drive and a flick off the pads are separated by little more
than the angle of the bat face at contact, which is frequently not resolvable in a single
frame. Several \texttt{legglance-flick} samples would be read as drives by eye. The
second closest pair is \texttt{pullshot} and \texttt{sweep}, which share a horizontal
bat path but are separated by a clear difference in body height --- the sweep played from
a lowered or kneeling position, the pull from a standing one --- making them
comparatively easy to tell apart.
Separately, \texttt{legglance-flick} images are roughly three times larger in both
median area and mean file size than the other three classes, which is a potential
shortcut feature the network could exploit in place of posture.

\textbf{Finding 4 --- Geometry, format and channel inconsistency.} Images span
168--1280\,px in width and 157--1280\,px in height across 84 distinct size combinations,
with a median of $275 \times 191$. Formats are mixed: 4{,}586 PNG, 137 JPEG and 1 WebP,
so a minority of images carry JPEG compression artefacts the majority do not. Colour
modes are likewise mixed: 4{,}586 images are RGBA and 138 are RGB. No files were corrupt
or unreadable. \emph{Consequence:} every image is explicitly converted to three-channel
RGB and resized to $224 \times 224$ before training. The RGBA conversion is mandatory,
not cosmetic --- ImageNet-pretrained backbones accept three channels only.

\subsection{Preprocessing and augmentation plan}
\label{sec:prep}

\begin{tabularx}{\linewidth}{@{}p{4.0cm}p{4.0cm}X@{}}
\toprule
\textbf{Decision} & \textbf{Choice} & \textbf{Driven by} \\
\midrule
Channel conversion & \textbf{Force RGB} & Finding 4 --- 97\% of images are RGBA \\
Resize & $224 \times 224$ & Finding 4 --- 84 distinct sizes; backbone input \\
Normalisation & ImageNet mean/std & Transfer learning in Submission 2 \\
Horizontal flip & \textbf{Not applied} & Finding 3 --- see note below \\
Rotation & $\pm 10^{\circ}$ & Finding 4 --- realistic broadcast camera tilt \\
Brightness / contrast & $\pm 20\%$ & Finding 4 --- varying stadium illumination \\
Random resized crop & scale 0.8--1.0 & Finding 3 --- forces reliance on the batsman \\
Class weights & \textbf{Not applied} & Finding 1 --- ratio is 1.12\,:\,1 \\
Split strategy & Grouped, stratified, 70/15/15 & Finding 2 --- duplicate-group leakage \\
\bottomrule
\end{tabularx}

In code, the channel conversion is a single call that must not be omitted:

\begin{quote}
\texttt{Image.open(path).convert("RGB")}
\end{quote}

\begin{warnbox}
\textbf{On horizontal flipping.} This is the augmentation most commonly applied by
default in image classification pipelines, and it is incorrect for this dataset. A sweep
and a leg glance are defined in part by which side of the body the ball is played to.
Mirroring an image therefore changes the correct label, injecting systematic label noise
rather than useful invariance. It is excluded deliberately.
\end{warnbox}

\clearpage
\begin{landscape}
\section{Related work}

Peer-reviewed papers and preprints only. Note that the great majority of prior work on
cricket shot classification operates on \emph{video}; this is the gap the present project
addresses.

\footnotesize
\setlength{\LTleft}{\fill}
\setlength{\LTright}{\fill}
\renewcommand{\arraystretch}{1.15}
\begin{longtable}{@{}p{0.5cm}p{3.7cm}p{3.0cm}p{3.7cm}p{5.8cm}p{5.0cm}@{}}
\toprule
\textbf{\#} & \textbf{Paper (author, year)} & \textbf{Dataset} & \textbf{Model(s)} & \textbf{Approach} & \textbf{Result} \\
\midrule
\endfirsthead
\toprule
\textbf{\#} & \textbf{Paper} & \textbf{Dataset} & \textbf{Model(s)} & \textbf{Approach} & \textbf{Result} \\
\midrule
\endhead
\bottomrule
\endfoot

1 & Sen, Deb, Dhar \& Koshiba (2021), \textit{Sensors} --- ``CricShotClassify''~\cite{sen2021}
  & CricShot10 (own; 10 classes, video)
  & Conventional CNN, dilated CNN, VGG16 / InceptionV3 / Xception / DenseNet169, each $+$ GRU
  & Transfer learning with frozen backbones, then selective unfreezing of the final 4 and final 8 VGG16 layers; GRU for temporal dependency
  & 86\% (fully frozen VGG16--GRU); 93\% (final 4 or 8 layers fine-tuned) \\
\addlinespace

2 & Kang (2025), \textit{arXiv:2510.09187}~\cite{kang2025}
  & CricShot10
  & Seven models: CNN$+$GRU, Bi-LSTM $+$ attention, ViT, transfer learning, EfficientNet-B0 $+$ GRU
  & Unified re-implementation of prior work under a single standardised evaluation protocol
  & 92.25\% (proposed). \textbf{Re-implementations of three prior published models scored 46.0\%, 55.6\% and 57.7\% against originally reported 96\%, 99.2\% and 93\%} \\
\addlinespace

3 & Dey \& Biswas (2024), \textit{Int.\ J.\ Engineering} --- ``Shot-ViT''~\cite{dey2024}
  & CBSId (own; 2{,}160 annotated \textbf{images}, 7 classes)
  & Fine-tuned Vision Transformer, compared against VGG19, ResNet50, I-AlexNet, ViT\_B32
  & Image-level rather than video-level classification; ViT fine-tuned for cricket batting shots
  & Shot-ViT outperforms all four compared CNN and ViT baselines \\
\addlinespace

4 & Kumar et al.\ (2019), \textit{IJRTE}~\cite{kumar2019}
  & Own video dataset
  & Modified LRCN (CNN $+$ LSTM)
  & CNN for spatial features, LSTM for temporal modelling
  & 96\%, but on a reduced set of \textbf{5} shot classes rather than the full 10 \\
\addlinespace

5 & Bhat et al.\ (2024), \textit{IJERCSE}~\cite{bhat2024}
  & Cricket shot video dataset
  & CNN$+$RNN, attention network, ViT
  & Direct comparison of three deep learning paradigms on one dataset
  & 99.2\% (CNN$+$RNN), 99.19\% (attention), 98.9\% (ViT) --- not reproduced by the independent re-implementation in row 2 \\
\addlinespace

6 & Islam, Hassan \& Khan (2019), \textit{arXiv:1909.01228}~\cite{islam2019}
  & BowlersNet (own; 8{,}100 \textbf{images}, 18 bowlers)
  & VGG16 $+$ transfer learning
  & Single-frame image classification of cricket bowling actions; selective layer freezing
  & 93.3\% average test accuracy \\

\end{longtable}
\renewcommand{\arraystretch}{1.25}
\normalsize
\end{landscape}

\subsection{What this survey establishes}

\textbf{1. The field is overwhelmingly video-based.} Rows 1, 2, 4 and 5 all model
temporal sequences. Only rows 3 and 6 operate on single images, and row 6 classifies
bowlers rather than batting shots. Single-frame batting shot classification is therefore
genuinely under-examined, which is the niche this project occupies.

\textbf{2. Published accuracies in this area are unreliable.} Kang~\cite{kang2025}
re-implemented prior models under a standardised protocol and recovered roughly half the
originally reported accuracy in each case. Inconsistent data partitioning is a plausible
contributor --- precisely the failure mode that the duplicate audit in Section~\ref{sec:dupaudit} is
designed to eliminate here.

\begin{keybox}
This is why the leakage control in this project is not incidental tidiness. It is a
direct response to a documented reproducibility problem in the immediate literature, and
it should be cited when presenting this project's own results.
\end{keybox}

{\footnotesize\color{muted}\textit{Note:} citations in rows 4 and 5 of the table above
are reported via the survey in \cite{kang2025}; the original publications are not
openly indexed, so their figures are as reported there.}

\textbf{3. Attention appears, but is never isolated.} Row 5 reports an attention-based
network, but introduces it alongside other architectural changes, so the contribution of
attention alone cannot be separated from the rest. Submission 4 of this project holds
architecture, data, split, seed and all hyperparameters constant and varies only the
presence of the attention module, making the measured difference attributable.

% ========================================================================
\section{Protocol for subsequent submissions}

The following are fixed here and held constant throughout, so that every later
comparison is attributable.

\begin{tabularx}{\linewidth}{@{}>{\bfseries}p{5.2cm}X@{}}
\toprule
Setting & Value \\
\midrule
Split & 70 / 15 / 15, grouped by duplicate identity, stratified by class \\
Random seed & 42 \\
Input resolution & $224 \times 224$ \\
Normalisation & ImageNet channel statistics \\
Reported metrics & Accuracy, precision, recall, F1 (per class and macro), confusion matrix, training and testing time \\
Split source & \texttt{results-data/split\_manifest.csv} --- loaded, never regenerated \\
\bottomrule
\end{tabularx}

\subsection*{Planned work}

\begin{tabularx}{\linewidth}{@{}p{1.0cm}p{5.4cm}X@{}}
\toprule
\textbf{\#} & \textbf{Deliverable} & \textbf{Scope} \\
\midrule
2 & Pretrained models & Three ImageNet backbones (ResNet50, EfficientNet-B0, DenseNet121), each evaluated frozen and fine-tuned, 30--50 epochs --- six training runs \\
3 & Custom CNN from scratch & Four convolutional blocks with batch normalisation and ReLU, global average pooling, fully connected head; augmentation and early stopping \\
4 & CNN with attention & The Submission 3 architecture with a CBAM module inserted after each convolutional block. Nothing else changes. Grad-CAM comparison of where each model attends \\
\bottomrule
\end{tabularx}

% ========================================================================
\section{Submission checklist}

\begin{enumerate}[label=$\square$\ ,leftmargin=2em]
\item Dataset added as a notebook input; \texttt{DATA\_DIR} resolves without error
\item Notebook executed with \textbf{Restart \& Run All}, so all plots and tables are stored in the \texttt{.ipynb}
\item Licence and year filled in from the Kaggle page
\item Confusion pair in Finding 3 confirmed against the Section 2.2 sample grid
\item \texttt{results-data/split\_manifest.csv} generated and committed
\item Reference links verified against the original publications
\item Notebook pushed to GitHub with outputs visible; repository link submitted
\end{enumerate}

% ========================================================================
\begin{thebibliography}{9}
\small

\bibitem{sen2021}
A.~Sen, K.~Deb, P.~K.~Dhar, and T.~Koshiba,
``CricShotClassify: An Approach to Classifying Batting Shots from Cricket Videos Using a
Convolutional Neural Network and Gated Recurrent Unit,''
\textit{Sensors}, 2021.
\url{https://pmc.ncbi.nlm.nih.gov/articles/PMC8072636}

\bibitem{kang2025}
S.~Kang,
``Modern Deep Learning Approaches for Cricket Shot Classification: A Comprehensive
Baseline Study,''
\textit{arXiv:2510.09187}, 2025.
\url{https://arxiv.org/abs/2510.09187}

\bibitem{dey2024}
A.~Dey and S.~Biswas,
``Shot-ViT: Cricket Batting Shots Classification with Vision Transformer Network,''
\textit{International Journal of Engineering}, 2024.
\url{https://www.ije.ir/article_195819.html}

\bibitem{kumar2019}
B.~Kumar et al.,
``Video based cricket shot classification using modified LRCN approach,''
\textit{International Journal of Recent Technology and Engineering}, vol.~8, no.~3, 2019.

\bibitem{bhat2024}
I.~Bhat et al.,
``Classification of cricket shots from cricket videos using deep learning models,''
\textit{International Journal of Engineering Research in Computer Science and
Engineering}, vol.~11, no.~5, 2024.

\bibitem{islam2019}
M.~N.~A.~Islam, T.~B.~Hassan, and S.~K.~Khan,
``A CNN-based approach to classify cricket bowlers based on their bowling actions,''
\textit{arXiv:1909.01228}, 2019.
\url{https://arxiv.org/abs/1909.01228}

\end{thebibliography}


\end{document}
